%%%%%%%%%%%%%%%%%%%%%%%%%%%%%%%%%%%%%%%%%%%%%%%%%%%%%%%%%%%%%%%%%%%%%%%%%%%%%%
%%%%%                                                                    %%%%%
%%%%% Sample Exam using accessibleexam.cls                               %%%%%
%%%%% Includes math and tikz diagram as examples.                        %%%%%
%%%%% Author: Kris Hollingsworth, PhD                                    %%%%%
%%%%% Original Date: 4/25/2026                                           %%%%%
%%%%%                                                                    %%%%%
%%%%%%%%%%%%%%%%%%%%%%%%%%%%%%%%%%%%%%%%%%%%%%%%%%%%%%%%%%%%%%%%%%%%%%%%%%%%%%
\DocumentMetadata{
  tagging=on,
  pdfstandard=ua-2,
  lang=en-US,
  testphase={phase-III,math,table}
}
\documentclass{accessibleexam}
%%%%%%%%%%%%%%%%%%%% Required for PDF File Metadata.
\usepackage[
  pdfauthor={Kris Hollingsworth, PhD},
  pdftitle={Exam X: Title of Exam}
]{hyperref}
%%%%%%%%%%%%%%%%%%%% Document Macros
\course{MATHXXX: The Best Math Course}
\exam{Exam X: The Best Math Exam}
\examdateval{February 30th, 2026}
\setinstructions{The best exam instructions go here.}
%=======================================================================================================================
%%% Load Custom Packages
\RequirePackage{tikz, pgfplots}
\RequirePackage{array}
\pgfplotsset{compat=1.18}
\usepgfplotslibrary{fillbetween}
%=======================================================================================================================
\begin{document}
%=======================================================================================================================
%%%%%%%%%%%%%%%%%%%%%%%%
%%%%%%%%%%%%% Cover Page
%%%%%%%%%%%%%%%%%%%%%%%%
%%% Give points per page as comma separated list. Will build grade table and automatically place point total per page in footer.
\makeexamcover{5,10,20}
%=======================================================================================================================
%%%%%%%%%%%%%%%%%%%%%%%%
%%%%%%%%%%%%%%%%% Page 1
%%%%%%%%%%%%%%%%%%%%%%%%
\question[2]{Please compute the following sums.}

\begin{enumerate}
  \item 
  $2+2 = $
\vfill
  \item 
  $\pi + e = $
\vfill
\end{enumerate}

\question[3]{Define a linear function.}
\vfill
%=======================================================================================================================
%%%%%%%%%%%%%%%%%%%%%%%%
%%%%%%%%%%%%%%%%% Page 2
%%%%%%%%%%%%%%%%%%%%%%%%
\newpage
\question[10]{All your base are belong to us.}

%=======================================================================================================================
%%%%%%%%%%%%%%%%%%%%%%%%
%%%%%%%%%%%%%%%%% Page 3
%%%%%%%%%%%%%%%%%%%%%%%%
\newpage

\question[20]{Consider the following linear function.}

{\centering
\begin{tikzpicture}[alt={The best linear function.}]
\begin{axis}[
    xmin=-.5,xmax=3.5,
    ymin=-3.5,ymax=3.5,
    grid=major,
    xtick={-10,-9,...,10},
    ytick={-5,-4,...,5},
    grid style={line width=.1pt, draw=gray!10},
    major grid style={line width=.2pt,draw=gray!20},
    axis lines=middle,
    enlargelimits={abs=0.5},
    legend style={at={(.85,1)}, anchor=north west},
    axis line style={latex-latex},
  height=.3\textwidth,
  width=.8\textwidth
]
\addplot [
    name path=line,
    very thick,
    domain=0:3, 
    samples=2, 
    color=blue,
]
{x+1} node[above,pos=.5,rotate=15] {y=x+1};
\addlegendentry{\(x+1\)}
\end{axis}
\end{tikzpicture}\par}

\begin{enumerate}
  \item What is the slope of the linear function?
  \vfill
  \item What is the $y$-intercept of the function?
  \vfill
  \item What is the $x$-intercept of the function?
  \vfill
  \vfill
\end{enumerate}
\end{document}
